\documentclass[12pt]{article}

\usepackage[utf8]{inputenc}
\usepackage[romanian]{babel}
\usepackage{amsmath, amssymb, amsthm}
\usepackage{geometry}
\geometry{a4paper, margin=2.5cm}

\title{Analiză Funcțională — Spațiu Reflexiv}
\author{}
\date{}

\theoremstyle{definition}
\newtheorem{definition}{Definiție}

\begin{document}

\maketitle

\begin{definition}
Fie $(E,\|\cdot\|)$ un spațiu normat. Dualul său este


\[
E^* = \{ f : E \to \mathbb{R} \text{ liniar și continuu} \}.
\]


Dualul dualului este


\[
E^{**} = (E^*)^*.
\]


\end{definition}

\begin{definition}
Aplicația canonică (naturală)


\[
J : E \to E^{**}, \qquad J(x)(f) = f(x),\quad f \in E^*,
\]


se numește \emph{imbedding-ul canonic}.
\end{definition}

\begin{definition}
Un spațiu normat $E$ se numește \emph{reflexiv} dacă aplicația canonică


\[
J : E \to E^{**}
\]


este surjectivă, adică


\[
J(E) = E^{**}.
\]


\end{definition}

\end{document}
